\section{Avoiding exceptional states and the mixing scale}
\label{rcut:sec:dynamics}

The gradient estimate contracts outside an exceptional set of small
Gibbs mass.  To use it from an arbitrary initial configuration, we must
also control the time the chain spends in that set.  We first obtain
this control after a constant burn-in, then propagate gradients over
logarithmic times.  The resulting concentration bounds will locate the
median mixing time on that same scale.

Fix $\beta\in[0,1/2)$.  Intersect the
good-disorder events in Propositions~\ref{rcut:prop:good-disorder},
\ref{rcut:prop:uniform-gap}, \ref{rcut:prop:entropy}, and
\ref{rcut:prop:typical-gradient}, and denote the resulting event by
$\mathcal G_n$.  Then $\P_J(\mathcal G_n)\to1$.
We collect the inputs on this event before starting the argument.
There are constants $K,C_E,C_L,\gamma,\kappa>0$ and $A\ge0$, depending
only on the fixed $\beta$, such that the following bounds hold for all
sufficiently large $n$ and every realization in $\mathcal G_n$.
The interaction satisfies $\|J\|\le K$, the continuous-time gap is at
least $\gamma$, and the transition density has the entropy bound
\begin{equation}
 \Ent_\mu\!\left(\frac{S_t(x,\cdot)}{\mu}\right)
 \le C_E n e^{-t/C_L}
 \qquad(x\in\Omega_n,\ t\ge0).
 \label{rcut:eq:dyn-entropy-input}
\end{equation}
Here the quotient denotes the density with respect to $\mu$.
There is also a set $\mathcal B=\mathcal B(J)\subseteq\Omega_n$ such that
\begin{equation}
 \mu(\mathcal B)\le e^{-\kappa n},\qquad
 (\partial_r-L)U(S_r f)
 \le \bigl(-\kappa+A\1_{\mathcal B}\bigr)U(S_r f)
 \quad(r\ge0)
 \label{rcut:eq:dyn-gradient-input}
\end{equation}
for every real function $f$ on $\Omega_n$.  The same $\kappa$ can be
used for the exceptional mass and the contraction rate by shrinking it
and enlarging $A$.  From now on we work deterministically on
$\mathcal G_n$: both the constants and the sufficiently large dimension
thresholds are uniform over this event.  Suppressing the fixed disorder,
we write $\P_x,\E_x$ for the law and expectation of the continuous-time
chain started at $x$.

\subsection{Gradient propagation and jump concentration}

We begin by expressing the gradient bound along a chain trajectory.
The exponential weight below records contraction outside the exceptional
set and the possible loss during visits to it.  A rough bound on this
weight will suffice for the initial burn-in.

\begin{lemma}[Feynman--Kac comparison]
\label{rcut:lem:feynman-kac}
Put $V=-\kappa+A\1_{\mathcal B}$.  For every real $f$, every
$0\le s\le d$, and every initial configuration $x$,
\begin{equation}
 U(S_d f)(x)
 \le \E_x\!\left[
    \exp\!\left(\int_0^s V(X_u)\,du\right)
    U(S_{d-s}f)(X_s)
 \right].
 \label{rcut:eq:dyn-feynman-kac}
\end{equation}
Consequently,
\begin{equation}
 \|U(S_r f)\|_\infty
 \le e^{Ar}\|U(f)\|_\infty
 \qquad(r\ge0).
 \label{rcut:eq:dyn-rough-envelope}
\end{equation}
\end{lemma}

\begin{proof}
Let $G(r,x)=U(S_r f)(x)$.  Equation~\eqref{rcut:eq:dyn-gradient-input}
gives $-\partial_rG+L G+VG\ge0$.  Following the chain forward while
running the semigroup parameter backward therefore gives nonnegative
drift for
\[
 \exp\!\left(\int_0^u V(X_v)\,dv\right)G(d-u,X_u),
 \qquad 0\le u\le s,
\]
so the finite state space and bounded time interval justify taking
expectations.  This proves Equation~\eqref{rcut:eq:dyn-feynman-kac}.
For the rough envelope, take $s=d=r$ and bound the exponential weight
using $V\le A$.  This gives Equation~\eqref{rcut:eq:dyn-rough-envelope}.
\end{proof}

To convert gradient envelopes into fluctuation bounds we use
Lemma~\ref{lem:input-mgf}, with $U(f)=\|df\|^2$.  Its complete
finite-chain proof uses only the heat-bath jump rates, so it applies
throughout the present argument without a temperature assumption.

\subsection{Uniform avoidance after a constant burn-in}

Small Gibbs mass alone does not control a chain started inside the
exceptional set.  We first show that entropy dissipation moves a fixed
amount of probability a macroscopic Hamming distance away from that
set.  Concentration of the distance then makes the probability of a
visit at any later deterministic time very small.

\begin{proposition}
\label{rcut:prop:bad-avoidance}
There are constants $d_0,c_{\mathrm a}>0$, independent of $n$ and
the realization in $\mathcal G_n$, such that, for all sufficiently
large $n$,
\begin{equation}
 \sup_{x\in\Omega_n}\P_x(X_u\in\mathcal B)
 \le q_n:=e^{-c_{\mathrm a}\sqrt n}
 \qquad(u\ge d_0).
 \label{rcut:eq:dyn-bad-avoidance}
\end{equation}
\end{proposition}

\begin{proof}
If $\mathcal B$ is empty there is nothing to prove.  Otherwise we
first enlarge it by a small fixed fraction of the coordinates and show
that the enlargement still has exponentially small Gibbs mass.  Let
$d_H$ denote Hamming distance and, for $0<\delta<1/2$, set
\[
 \mathcal B^{(\delta)}
 =\{x:d_H(x,\mathcal B)\le\lfloor\delta n\rfloor\}.
\]
If $d_H(x,y)\le\delta n$, then
\[
 \left|\frac{x^{\mathsf T}Jx-y^{\mathsf T}Jy}{2}\right|
 =\frac12\left|(x-y)^{\mathsf T}J(x+y)\right|
 \le2K\sqrt\delta\,n.
\]
Thus the energy change within the neighborhood is controlled by the
operator norm of the interaction.  Counting the neighbors of each
$y\in\mathcal B$, with multiplicity, gives
\begin{equation}
 \mu(\mathcal B^{(\delta)})
 \le \mu(\mathcal B)e^{2K\sqrt\delta\,n}
       \sum_{j\le\delta n}\binom nj
 \le \exp\!\left[-n\bigl(\kappa-2K\sqrt\delta-h(\delta)\bigr)\right],
 \label{rcut:eq:dyn-neighborhood}
\end{equation}
where $h(\delta)=-\delta\log\delta-(1-\delta)\log(1-\delta)$.
The binomial bound follows by inserting
$z=\delta/(1-\delta)<1$ into $(1+z)^n$ and using
$z^j\ge z^{\delta n}$ for $j\le\delta n$.
Choose a fixed $\delta$ so small that
$2K\sqrt\delta+h(\delta)\le\kappa/2$.  Equation~\eqref{rcut:eq:dyn-neighborhood}
then bounds the neighborhood mass by $e^{-\kappa n/2}$.

We next use this small equilibrium mass to bound the probability of
the neighborhood after burn-in.  For probability measures $\nu,\mu$
and an event of probabilities $p,q$, respectively, conditioning on the
event and its complement gives
\[
 \Ent_\mu\!\left(\frac{\nu}{\mu}\right)
 \ge p\log\frac pq+(1-p)\log\frac{1-p}{1-q}
 \ge p\log\frac1q-\log2.
\]
Apply this inequality to $\nu=S_{d_0}(x,\cdot)$ and the neighborhood,
choosing the burn-in time to be
\[
 d_0=1+C_L\log\max\{1,8C_E/\kappa\}.
\]
By Equation~\eqref{rcut:eq:dyn-entropy-input}, for sufficiently large $n$,
\[
 \P_x(X_{d_0}\in\mathcal B^{(\delta)})
 \le\frac{2C_E}{\kappa}e^{-d_0/C_L}
       +\frac{2\log2}{\kappa n}
 \le\frac12.
\]

At least half the law is now outside the neighborhood.  To strengthen
this fact to avoidance of the original set, consider
$g(x)=d_H(x,\mathcal B)/\sqrt n$.  Its half-differences have magnitude
at most $1/(2\sqrt n)$, so $U(g)\le1/4$.
Outside $\mathcal B^{(\delta)}$ the distance is greater than $\delta n$.
The integer rounding in the neighborhood definition therefore gives
\[
 S_{d_0}g(x)\ge\frac{\delta\sqrt n}{2}.
\]
Equation~\eqref{rcut:eq:dyn-rough-envelope} permits the envelope
$b(r)=e^{Ar}/4$ for $0\le r\le d_0$.
Put $Q=e^{Ad_0/2}$.  Then $B\le Q^2/4$ and
$I\le d_0Q^2/4$ in Lemma~\ref{lem:input-mgf}.
Taking $\theta=-1/Q$ in Lemma~\ref{lem:input-mgf} yields
\[
 \P_x(X_{d_0}\in\mathcal B)
 \le\exp\!\left(\frac{e d_0}{2}
                 -\frac{\delta\sqrt n}{2Q}\right)
 \le e^{-c_{\mathrm a}\sqrt n},
 \qquad c_{\mathrm a}=\frac{\delta}{4Q},
\]
after increasing the fixed minimum dimension.  The event on the left
is exactly $\{g(X_{d_0})=0\}$, so the negative exponential moment
controls precisely the required lower tail.  This proves avoidance at
the burn-in time.  To obtain the same bound at every deterministic
$u\ge d_0$, apply the Markov property to the last burn-in interval:
\[
 \P_x(X_u\in\mathcal B)
 =\E_x\!\left[\P_{X_{u-d_0}}(X_{d_0}\in\mathcal B)\right]
 \le q_n.
\]
Every choice depended only on the constants collected at the start of
this section.  The estimate is therefore uniform in the starting point
and in the disorder on the good event.
\end{proof}

The pointwise-in-time avoidance bound also controls occupation time by
integration.  We now return to the Feynman--Kac weight: on a logarithmic
time horizon, an exceptional occupation larger than one has such small
probability that even the rough exponential envelope is harmless.

\begin{lemma}[Propagation on a logarithmic time horizon]
\label{rcut:lem:dyn-occupation-propagation}
Fix $D_*>0$ and set
\begin{equation}
 T_n=D_*\log n,\qquad
 C_{\mathrm{dyn}}=e^{A(d_0+1)},\qquad
 r_n=T_ne^{AT_n}q_n
     =D_*\log n\,n^{AD_*}e^{-c_{\mathrm a}\sqrt n}.
 \label{rcut:eq:dyn-occupation-error}
\end{equation}
For every $0\le s\le d\le T_n$, every real $f$, and every $x$,
\begin{equation}
 \begin{split}
 U(S_d f)(x)
 &\le C_{\mathrm{dyn}}e^{-\kappa s}
      S_s\!\left[U(S_{d-s}f)\right](x)\\
 &\quad+r_n\|U(S_{d-s}f)\|_\infty.
 \end{split}
 \label{rcut:eq:gradient-propagation}
\end{equation}
For every fixed $p>0$, $r_n=o(n^{-p})$.
\end{lemma}

\begin{proof}
Measure only the exceptional occupation after burn-in: let
$Z_n=\int_{d_0}^{T_n}\1_{\mathcal B}(X_u)\,du$ when
$T_n\ge d_0$, and put $Z_n=0$ otherwise.  Integrating
Proposition~\ref{rcut:prop:bad-avoidance} and applying Markov's
inequality gives
\[
 \P_x(Z_n>1)\le\E_x Z_n\le T_nq_n.
\]
On $\{Z_n\le1\}$, the total bad-set occupation up to any
$s\le T_n$ is at most $d_0+1$.  Therefore
\[
 \exp\!\left(\int_0^s V(X_u)\,du\right)
 \le C_{\mathrm{dyn}}e^{-\kappa s}
 \quad\hbox{on }\{Z_n\le1\}.
\]
The rough bound $e^{AT_n}$ applies to the exponential weight on the
complementary event.  Split the expectation in
Lemma~\ref{rcut:lem:feynman-kac} accordingly.  For occupation at most one, use the displayed weight bound and enlarge
the expectation of the nonnegative endpoint quantity to the whole
space.  For larger occupation, use its supremum and the probability
bound above.  These two contributions are exactly the terms in
Equation~\eqref{rcut:eq:gradient-propagation}; occupation and endpoint
gradient need not be independent.  This also covers $s<d_0$, since the
occupation up to $s$ is then at most $d_0$.  Finally,
$e^{-c_{\mathrm a}\sqrt n}$ dominates every fixed power of $n$, which
proves the stated decay of the error.
\end{proof}

\subsection{Linear observables and the order of the mixing time}

The propagated gradient has a bounded time integral for every linear
observable with bounded coefficient norm.  The jump exponential-moment
estimate therefore gives fluctuations of bounded size, uniformly over
all such observables and all starts.  We will use this uniformity to
choose a distinguishing observable after the disorder and the initial
configuration have been selected.

\begin{proposition}[Uniform concentration for linear observables]
\label{rcut:prop:linear-concentration}
Fix $D_*>0$ and $L\ge0$.  For $a\in\R^n$ write
$\ell_a(x)=a^{\mathsf T}x$, and define
\[
 B_L=L^2(C_{\mathrm{dyn}}+1),\qquad
 I_L=L^2\left(\frac{C_{\mathrm{dyn}}}{\kappa}+1\right).
\]
For all sufficiently large $n$, every $\|a\|\le L$, every $x$,
every $0\le d\le D_*\log n$, and every $\theta\in\R$,
\begin{equation}
 \E_x\exp\!\left(\theta\bigl(\ell_a(X_d)-S_d\ell_a(x)\bigr)\right)
 \le M_{\theta,L}
 :=\exp\!\left(2\theta^2e^{2|\theta|\sqrt{B_L}}I_L\right),
 \label{rcut:eq:dyn-linear-mgf}
\end{equation}
and
\begin{equation}
 \Var_x(\ell_a(X_d))\le4I_L,\qquad
 \Var_\mu(\ell_a)\le\gamma^{-1}L^2,\qquad
 \mu(\ell_a)=0.
 \label{rcut:eq:dyn-linear-variance}
\end{equation}
These bounds hold simultaneously over all the indicated vectors and
starting points, and uniformly over $\mathcal G_n$.
\end{proposition}

\begin{proof}
For a linear observable, $U(\ell_a)=\|a\|^2$, so the endpoint gradient
in the propagation estimate is constant.  Taking $s=d=r$ in
Equation~\eqref{rcut:eq:gradient-propagation} gives the deterministic
envelope
\[
 \|U(S_r\ell_a)\|_\infty
 \le b(r):=L^2\bigl(C_{\mathrm{dyn}}e^{-\kappa r}+r_n\bigr),
 \qquad 0\le r\le T_n.
\]
For sufficiently large $n$, both $r_n\le1$ and $T_nr_n\le1$.
Consequently $\sup b\le B_L$ and $\int_0^{T_n}b(r)\,dr\le I_L$.
These are the two quantities needed by Lemma~\ref{lem:input-mgf}, which
gives both the transient exponential-moment bound and the variance bound.

The equilibrium variance uses the gap input instead.  Applying its
Dirichlet-form inequality gives
\[
 \Var_\mu(\ell_a)
 \le\gamma^{-1}\mu\!\left[\ell_a(-L)\ell_a\right]
 =\gamma^{-1}\sum_i a_i^2\mu(v_i)
 \le\gamma^{-1}\|a\|^2.
\]
The energy identity follows by conditional expectation:
$\mu(x_i m_i)=\mu(m_i^2)$, so the coordinate energy is $\mu(v_i)$.
Spin inversion in the zero-field model supplies the zero mean in
Equation~\eqref{rcut:eq:dyn-linear-variance}.  Since the gradient envelope
holds on the entire coefficient ball, none of these bounds requires
$a$ to be chosen before the disorder or the starting configuration.
\end{proof}

We can now place the mixing time between logarithmic bounds.  For the
lower bound, a spin correlation will produce a linear observable whose
mean is still large while both variances remain bounded.  For the upper
bound, the entropy estimate already available at the start of the
section suffices.

\begin{proposition}[The median mixing scale]
\label{rcut:prop:mixing-scale}
There is a constant $C<\infty$ such that, uniformly over
$\mathcal G_n$ and all sufficiently large $n$,
\begin{equation}
 d_c\!\left(\frac14\log n\right)
 \ge1-Cn^{-1/2}.
 \label{rcut:eq:dyn-strong-lower}
\end{equation}
The continuous median mixing time
\[
 t_n^{\rm med}=\inf\{t\ge0:d_c(t)\le1/2\}
\]
satisfies $d_c(t_n^{\rm med})=1/2$ and, for some constants $0<c<C<\infty$,
\begin{equation}
 c\log n\le t_n^{\rm med}\le C\log n.
 \label{rcut:eq:dyn-median-scale}
\end{equation}
\end{proposition}

\begin{proof}
Write $\sigma_i(x)=x_i$.  Spin inversion gives
$\mu(\sigma_i)=0$ and $\mu(\sigma_i^2)=1$, while
\[
 \mu\!\left[\sigma_i(-L)\sigma_i\right]=\mu(v_i)\le1.
\]
Use the spectral decomposition of the finite-dimensional
self-adjoint operator $-L$.  The squared coefficients of $\sigma_i$
form nonnegative weights summing to one, and the preceding energy is
their mean eigenvalue.  Jensen's inequality for $\lambda\mapsto
e^{-d\lambda}$ therefore yields
\begin{equation}
 \mu\!\left[\sigma_i S_d\sigma_i\right]
 \ge\exp\!\left(-d\mu[\sigma_i(-L)\sigma_i]\right)
 \ge e^{-d}.
 \label{rcut:eq:dyn-spin-correlation}
\end{equation}
Sum these correlations over $i$.  Since their sum is an average over
the initial state under $\mu$, some configuration $x=x(d,J)$ satisfies
\[
 S_d\ell_{x/\sqrt n}(x)
 =\frac1{\sqrt n}\sum_i x_i S_d\sigma_i(x)
 \ge\sqrt n e^{-d}.
\]
The coefficient vector $x/\sqrt n$ has norm one, so
Proposition~\ref{rcut:prop:linear-concentration}, with $D_*=1$ and $L=1$,
applies even to this choice.  Let $V_*=4I_1$.  The displayed
observable has transient variance at most $V_*$, equilibrium mean
zero, and equilibrium variance at most $\gamma^{-1}$.
Let its transient mean be $m\ge\sqrt n e^{-d}$.  The event that the
observable exceeds $m/2$ separates the transient law from equilibrium.
Applying Chebyshev's inequality under both laws gives, for
$0\le d\le\log n$,
\[
 d_c(d)\ge1-\frac{4(V_*+\gamma^{-1})}{m^2}
 \ge1-\frac{4(V_*+\gamma^{-1})e^{2d}}{n}.
\]
Taking $d=(\log n)/4$ proves Equation~\eqref{rcut:eq:dyn-strong-lower}.

For the upper bound, the entropy-to-total-variation estimate proved in
Proposition~\ref{rcut:prop:entropy}, applied to
Equation~\eqref{rcut:eq:dyn-entropy-input}, gives
\[
 d_c(t)^2\le\frac12 C_E n e^{-t/C_L}.
\]
Hence
$t_n^{\rm med}\le C_L\log(2C_E n)$ for sufficiently large $n$.
Together with Equation~\eqref{rcut:eq:dyn-strong-lower}, this proves
Equation~\eqref{rcut:eq:dyn-median-scale}.
It remains to justify the exact median identity.  The function $d_c$
is continuous, as a maximum of finitely many continuous transition-law
distances, and the entropy bound shows that it tends to zero.  Since
$d_c(0)=1-\min_x\mu(x)\ge1-2^{-n}>1/2$, its first crossing is
attained and satisfies $d_c(t_n^{\rm med})=1/2$.
\end{proof}
